\documentclass[11pt]{article}

% Common packages
\usepackage{amsmath}
\usepackage{amssymb}
\usepackage{amsthm}
\usepackage{geometry}
\usepackage{enumitem}
\usepackage{graphicx}
\usepackage{titlesec}
\usepackage{tikz}
\usepackage{pgfplots}
\pgfplotsset{compat=1.18}

% Page setup
\geometry{margin=0.5in}
\setlist[itemize]{topsep=2pt,itemsep=2pt,parsep=0pt}

% --- Load hyperref late, then bookmark AFTER hyperref ---
\usepackage[hidelinks]{hyperref}
\usepackage{bookmark}

% Short labels
\newcommand{\thm}{\underline{\textbf{Thm.}} }
\newcommand{\defn}{\underline{\textbf{Def.}} }
\newcommand{\prop}{\underline{\textbf{Prop.}} }
\newcommand{\ex}{\underline{\textbf{Ex.}} }
\newcommand{\claim}{\underline{\textbf{Claim:}} }

\newcommand{\Gen}{\textsf{Gen}}
\newcommand{\Enc}{\textsf{Enc}}
\newcommand{\Dec}{\textsf{Dec}}

% Title info
\title{MATH 4350 Notes - Number Theory and Cryptography}
\author{Michael Lee}

\begin{document}
\maketitle

\tableofcontents
\newpage

\section{Syllabus / How to Class Works}
\begin{itemize}
    \item See course website: \href{https://www.math.wustl.edu/~matkerr/4350/}{https://www.math.wustl.edu/~matkerr/4350/}
\end{itemize}

\section{Introduction}
\begin{itemize}
    \item \textbf{GOAL:} How do we disguise messages s.t. only the sender and receiver know what the message is even when the encryption strategy is known?
\end{itemize}
\subsection{Greatest Common Divisors}
\begin{itemize}
    \item Let \(\mathbb{N} = \{1, 2, 3, \dots\}\) be the set of natural numbers.
    \item Let \(\mathbb{Z} = \{0, \pm 1, \pm 2, \dots\}\) be the set of integers.
    \item \defn Greatest Common Divisor (GCD): \((a,b) :=\) largest element of \(S(a,b) := \{d \in \mathbb{N} | d \mid a, b\}\)
    \item Properties of the GCD (true, but not proved in this class):
    \begin{itemize}
        \item \(a|b \wedge b|c \Rightarrow a|c\)
        \item \(a|b \wedge b|a = a = \pm b\)
        \item if \(a|b\) and \(a|c\) and \(x, y \in \mathbb{Z}\), then \(a|(bx + cy)\)
    \end{itemize}
    \item Recall that induction exists lmao
    \item \ex Find the GCD of \((a,b) = (162, 72)\)
    \begin{itemize}
        \item Use GCD props:
        \item \((0, a) = a = (a,a)\)
        \item \((a, b) = (b,a) = (a, -b) \)
        \item \((a, b-ma) = (a, b) \) for some \(m \in \mathbb{Z}\)
        \item Recall long division: \(A = \{a-bk | k \in \mathbb{Z}\, a-bk \geq 0\} \subseteq \mathbb{N} \cup \{0\}\) has least element \(r\) if \(r \geq b\), then \(r-b \geq 0 \Rightarrow \) \(r \in A\). Thus \(a = bq + r\) for some unique \(q, r \in \mathbb{Z}\) and \(0 \leq r < b\).
        \item Going back to the example,
        \[
        162 = 72 \cdot 2 + 18 \Rightarrow (162, 72) = (18, 72) = (72, 18)
        \]
        \[
        72 = 18 \cdot 4 + 0 \Rightarrow (72, 18) = (18, 0) = 18
        \]
        Thus, \((162, 72) = 18\). Let's generalize this:
    \end{itemize}
    \item \defn Euclidean Algorithm:
    \begin{itemize}
        \item \(a = bq_1 + r_1, \, 0 \leq r_1 < b\)
        \item \(b = r_1q_2 + r_2, \, 0 \leq r_2 < r_1\)
        \item \(r_1 = r_2q_3 + r_3, \, 0 \leq r_3 < r_2\)
        \item \(\vdots\)
        \item \(r_{n-1} = r_nq_{n+1} + r_{n+1}\)
        \item Eventually, \(r_{n+1} = 0\), so \(r_{n-1} = r_nq_{n+1} + 0\)
        \item Thus, \((a,b) = r_n\)
        \item Eventually, we will prove this is a polynomial time algorithm.
    \end{itemize}
    \item \thm \(S := \{ax+by | x,y \in \mathbb{Z}, ax+by > 0\} = (a,b) = g\) (i.e. the GCD)
    \begin{proof}
        First, we must show \(S \subseteq g\). We know that \(a = g \alpha\) and \(b = g \beta\) since \(g|a,b\) by defn for some \(x,y \in \mathbb{Z}\). \\

        Thus, any \(ax+by = g(\alpha x + \beta y)\) is divisible by \(g\). \\

        Next, we must prove that \(g \subseteq S\). We know that \(g = (a,b) = r_n = \) the last non-zero remainder in the Euclidean Algorithm. We will show all non-zero \(r_j\)'s are in \(S\) via induction:

        \textbf{Base Case:} \(r_1 = 1 \cdot a + (-q_1) \cdot b \in S\)

        \textbf{Inductive Hypothesis:} Assume \(r_j = a \cdot x_j + b \cdot y_j \in S \quad \forall j \in [1, \cdots, k-1]\).

        \textbf{Inductive Step:} We know:
        \[
        r_{k-2} = r_{k-1}q_k + r_k \Rightarrow r_k = r_{k-2} + (-q_k) r_{k-1}
        \]
        \[
            = a \cdot x_{k-2} + b \cdot y_{k-2} + (-q_k) (a \cdot x_{k-1} + b \cdot y_{k-1}) = a (x_{k-2} - q_k x_{k-1}) + b (y_{k-2} - q_k y_{k-1}) \in S
        \]
    \end{proof}
\end{itemize}

\subsection{Fermat's Little Theorem}
\begin{itemize}
    \item \thm \(a^p \equiv a \mod p\) for all \(a \in \mathbb{Z}\) and prime \(p\)
    \begin{proof}
        \textbf{Claim:} 
    \end{proof}
\end{itemize}

\end{document}
